% from ega4-17-fr.tex
\oldpage[IV-4]{107}

de sous-préschémas affines de $S$, de présentation finie sur $S$, deux à deux disjoints, de réunion $S$, et ayant la propriété suivante: pour chaque $i\in I$, il existe deux morphismes finis de présentation finie et surjectifs
\[
  S_i''\xrightarrow{g_i}S_i'\xrightarrow{h_i}S_i,
\]
où $h_i$ est étale et $g_i$ plat et radiciel, et un $S$-morphisme $S_i''\to X$ (autrement dit une $S_i''$-section de $X_i''=X\times_S S_i''$).
\end{proposition}

Il y a un recouvrement fini $(U_i)_{1\leq i\leq n}$ de $S$ par des ouverts affines; pour tout $i$, l'intersection
\[
  W_i=U_i\cap\left(\bigcup_{j<i}U_j\right)
\]
est quasi-compacte (I, 2.7) donc il existe un sous-préschéma fermé $V_i$ de $U_i$ ayant pour espace sous-jacent $U_i-W_i$ et défini par un idéal de type fini de l'anneau de $U_i$; donc $V_i$ est un sous-préschéma affine de $S$ qui est de présentation finie sur $S$ (I, 6.2). Il est clair qu'on peut se borner à prouver la proposition en y remplaçant $S$ par $V_i$ et $X$ par $X\times_S V_i$. Autrement dit, on peut déjà supposer $S$ affine. D'autre part, $X$ est réunion d'ouverts affines $W_\alpha$ et les $f(W_\alpha)$ sont constructibles dans $S$ (I, 8.4) et forment un recouvrement de $S$; par suite (I, 9.9) il existe une sous-famille finie $(W_{\alpha_j})_{1\leq j\leq m}$ telle que les $f(W_{\alpha_j})$ forment déjà un recouvrement de $S$. Soit alors $X'$ le préschéma somme des sous-préschémas induits sur les ouverts $W_{\alpha_j}$ de $X$; il est immédiat qu'il suffit de prouver la proposition en remplaçant $X$ par $X'$ puisqu'un $S$-morphisme $S_i''\to X'$ donne par composition un $S$-morphisme $S_i''\to X'\to X$.

On peut donc supposer $X$ affine et $f$ de présentation finie. On peut alors (8.9.1) écrire $X$ sous la forme $X_0\times_{S_0}S$, où $S_0$ est noethérien, $X_0\to S_0$ un morphisme de type fini, qu'on peut supposer surjectif (8.10.5, (vi)). Si la proposition est prouvée pour ce morphisme, elle en résultera aussitôt pour $X$ par changement de base $S\to S_0$. On peut donc supposer de plus $S$ noethérien et $f$ de type fini.

Soient $s_h$ $(1\leq h\leq r)$ les points maximaux de $S$. Montrons qu'il suffit de prouver l'énoncé en remplaçant $S$ par un voisinage ouvert assez petit $U_h$ de chacun des $s_h$, $X$ par l'image réciproque de $U_h$ dans $X$. En effet, supposons la proposition établie dans ce cas, et raisonnons par récurrence noethérienne ($0_{\mathrm I}$, 2.2.2) en supposant l'énoncé établi pour tout sous-préschéma fermé de $S$ ayant un espace sous-jacent $\neq S$. On peut supposer les $U_h$ deux à deux sans point commun; si $T$ est un sous-préschéma fermé ayant pour espace sous-jacent $S-(\bigcup_h U_h)$, l'hypothèse de récurrence entraîne que l'énoncé est vrai pour $T$; comme il est vrai aussi pour chacun des $U_h$, il l'est évidemment pour $S$.

Comme on peut évidemment (en remplaçant $S$ par $S_{\mathrm{red}}$ et $X$ par $X\times_S S_{\mathrm{red}}$) supposer $S$ réduit, chacun des $\sh{O}_{S,s_h}$ est un corps $k(s_h)$. Supposons que l'on ait prouvé la proposition lorsque $S$ est le spectre d'un corps et $f$ est de type fini. Alors l'existence des $U_h$ résulte de la méthode de (8.1.2, $a$)) et de (8.8.2, (i) et (ii)), (8.10.5, (vi), (vii) et (x)), (11.2.6, (ii)) et (17.7.8, (ii)).

Supposons donc $S=\Spec(k)$, où $k$ est un corps, $X$ étant de type fini sur $k$ et $X\neq\varnothing$. Comme $X$ est noethérien, il existe dans $X$ un point fermé $x$ ($0_{\mathrm I}$, 2.1.3), donc $k(x)$ est une extension finie de $k$ (I, 6.4.2). Il y a par suite une extension finie séparable $k'$ de $k$ telle que $k''=k(x)$ soit extension finie radicielle de $k'$. On répond alors à la question en prenant $I$ réduit à un élément, $S_i'=\Spec(k')$, $S_i''=\Spec(k'')$.

